\documentclass[acmsmall,screen]{acmart}
\citestyle{acmauthoryear}

%% --- ACM metadata --------------------------------------------------------
\acmJournal{PACMPL}
\acmVolume{1}
\acmNumber{HASKELL}
\acmYear{2026}
\acmDOI{}

%% --- Packages ------------------------------------------------------------
\usepackage{listings}
\usepackage{amsmath}
\let\Bbbk\relax
\usepackage{amssymb}
\usepackage{stmaryrd}   % \llbracket, \rrbracket
\usepackage{mathtools}
\usepackage{booktabs}
\usepackage{microtype}
\usepackage{enumitem}

%% --- Fix headheight for fancyhdr -----------------------------------------
\setlength{\headheight}{18.67603pt}
\addtolength{\topmargin}{-5.67603pt}

%% --- Haskell listings style ----------------------------------------------
\lstdefinelanguage{Haskell}{
  keywords={data, newtype, type, class, instance, where, let, in, do, if,
            then, else, case, of, import, module, deriving, infixl, infixr,
            return, Composable, Pledge, PledgeResult, liftPledge, inspect,
            leftQuotient, rightQuotient, concatenation, conjunction,
            empty, universe, runPledge,
            Bot, Epsilon, Single, Star, Atom, Wildcard, Usage,
            GuardedRE, LTL, MTL, Interval,
            PTrue, PFalse, Pred, Term, Values, ValAt},
  keywordstyle=\color{blue!70!black}\bfseries,
  sensitive=true,
  comment=[l]{--},
  morecomment=[s]{\{-}{-\}},
  commentstyle=\color{gray}\itshape,
  stringstyle=\color{orange!80!black},
  morestring=[b]",
  literate=
    {·}{{$\cdot$}}1
    {⊓}{{$\sqcap$}}1
    {∖}{{$\setminus$}}1
    {∕}{{$/$}}1
    {ε}{{$\varepsilon$}}1
    {⊤}{{$\top$}}1
    {->}{{$\to$}}2
    {=>}{{$\Rightarrow$}}2
    {forall}{{$\forall$}}1
    {>>}{{$>\!\!>$}}2
    {>>=}{{$>\!\!>\!\!=$}}3,
}
\lstset{
  language=Haskell,
  basicstyle=\ttfamily\small,
  columns=flexible,
  keepspaces=true,
  frame=none,
  xleftmargin=1em,
  breaklines=true,
  breakatwhitespace=true,
}

%% --- Math macros ---------------------------------------------------------
\newcommand{\RE}{\mathsf{RE}}
\newcommand{\GRE}{\mathsf{GRE}}
\newcommand{\LTL}{\mathsf{LTL}}
\newcommand{\MTL}{\mathsf{MTL}}
\newcommand{\pre}{\mathit{pre}}
\newcommand{\post}{\mathit{post}}
\newcommand{\fut}{\mathit{fut}}
\newcommand{\ret}{\mathit{ret}}
\newcommand{\anything}{\mathtt{anything}}
\newcommand{\finally}[1]{\mathtt{finally}(#1)}
\newcommand{\previously}[1]{\mathtt{previously}(#1)}
\newcommand{\never}[1]{\mathtt{never}(#1)}
\newcommand{\noUntil}[2]{\mathtt{noUntil}(#1,\,#2)}
\newcommand{\derives}[2]{\partial_{#1}(#2)}
\newcommand{\norm}[1]{\lceil #1 \rceil}
\newcommand{\sem}[1]{\llbracket #1 \rrbracket}
\newcommand{\univ}{\neg\varnothing}
\newcommand{\lquot}{\mathbin{\setminus}}   % left-quotient infix
\newcommand{\rquot}{\mathbin{/}}           % right-quotient infix

\begin{document}

%% --- Title / Authors -----------------------------------------------------
\title{Effectful Computations with Future Conditions:\texorpdfstring{\\}{}
       A Monadic Approach to Temporal Specification}

\author{Author Name}
\affiliation{%
  \institution{Institution}
  \city{City}
  \country{Country}
}
\email{author@institution.edu}

\begin{abstract}
Pre- and post-conditions are the standard vocabulary for modular program specification, yet they are silent about obligations that span multiple calls.
In real Haskell code such obligations are ubiquitous: an acquired mutex must eventually be released.
An opened handle must eventually be closed.
A pre/post pair on the acquiring operation cannot express these properties,
because the balancing operation lies in an unknown future context, tracking it means
threading auxiliary state by hand through every intermediate bind.

We present \emph{Pledge}, a Haskell library in which such obligations are
ordinary values.
A \lstinline{Pledge m eff a} carries, beside its result, what must have held
before it ran, what it produces, and what must still hold after.
One bind formula propagates all three by \emph{quotient}: binding strips
from a pending obligation whatever the continuation actually delivers.
What is left over --- the \emph{residual} --- names precisely what the
program still owes.

The formula needs only six operations, abstracted as a \lstinline{Composable}
type class, and we prove the monad laws in Lean~4 from twelve algebraic
axioms over that class, identifying exactly the two side conditions under
which our regular-expression instance $\RE$ satisfies them --- conditions
the obvious way of writing an annotation quietly violates, silently
reporting a violated contract for a correct program.
A soundness theorem completes the picture: a discharged residual entails
that the trace a run emitted met every obligation its primitives named.
On top of $\RE$ we build two front-ends that read the same residual two
different ways: an $\LTL$ layer that turns \lstinline{Pledge} into a
compositional \emph{runtime monitor} whose state is the residual itself, and
a bounded $\MTL$ layer that adds discrete-time intervals, compiled down to
$\RE$ by finite unrolling.
A separate $\GRE$ instance pairs an $\RE$ with a Presburger arithmetic
predicate, discharged by Z3, and because that predicate is checked over
symbolic heap addresses rather than a concrete run, $\GRE$ gives
\lstinline{Pledge} a \emph{static verification} mode alongside its runtime
one.
Because the underlying monad is arbitrary, the residual describes the path
a program actually took and can be computed beside a real effect handler.
\lstinline{Pledge} tells you what is still owed, and where; making it enforce
that, prove its own instances' laws automatically, and report empirical cost
are the open problems we lay out as a concrete roadmap.
\end{abstract}

\keywords{temporal specification, monadic effects, regular expressions,
          Brzozowski derivatives, left-quotient, right-quotient,
          Presburger arithmetic, metric temporal logic, runtime monitoring,
          static verification, Haskell, Lean~4}

\begin{CCSXML}
<ccs2012>
  <concept>
    <concept_id>10011007.10011006.10011008.10011009.10011012</concept_id>
    <concept_desc>Software and its engineering~Functional languages</concept_desc>
    <concept_significance>500</concept_significance>
  </concept>
  <concept>
    <concept_id>10011007.10011006.10011008.10011024.10011028</concept_id>
    <concept_desc>Software and its engineering~Data types and structures</concept_desc>
    <concept_significance>300</concept_significance>
  </concept>
  <concept>
    <concept_id>10011007.10011006.10011041</concept_id>
    <concept_desc>Software and its engineering~Software verification and validation</concept_desc>
    <concept_significance>300</concept_significance>
  </concept>
</ccs2012>
\end{CCSXML}
\ccsdesc[500]{Software and its engineering~Functional languages}
\ccsdesc[300]{Software and its engineering~Data types and structures}
\ccsdesc[300]{Software and its engineering~Software verification and validation}

\maketitle

%% =========================================================================
\section{Introduction}
\label{sec:why}
%% =========================================================================

Pre- and post-conditions are the workhorse of modular program specification.
A pre-condition constrains the state in which an operation may be invoked;
a post-condition describes what holds immediately after it returns.
Yet many correctness properties span a \emph{sequence} of operations, with
arbitrary computation in between.

Consider three scenarios from the Haskell ecosystem:
\emph{(i) Mutex discipline.}
An \lstinline{acquire(m)} must eventually be followed by
\lstinline{release(m)} on every code path.
A post-condition on \lstinline{acquire} records the acquired state, but the
\lstinline{release} that discharges it lies outside the scope that
post-condition describes.
\emph{(ii) Database savepoints.}
A \lstinline{beginTx} must eventually be committed or rolled back; a
pre/post pair on the call does not carry that obligation forward through the
intermediate code, so it must be re-stated by hand at every step in between.
\emph{(iii) File handles and heap cells.}
Opening a file in read-mode must be followed by a close; a heap cell
returned by \lstinline{malloc} must eventually be freed, and not freed
twice.

\paragraph{Existing remedies and their limits.}
The \lstinline{bracket} function lexically scopes cleanup but cannot
propagate obligations beyond the lexical scope.
\lstinline{ResourceT}~\cite{snoyman2011resourcet} threads cleanup through
every bind but models only the cleanup half of an obligation, discharging
it uniformly at scope exit.
Linear Haskell~\cite{bernardy2018linear} enforces ``consume exactly once'',
which does force discharge along every branch, but the obligation is
structural: it fixes \emph{that} a resource is consumed, not the order of the
operations leading there, and carries no quantitative reading.
Parameterised and indexed monads~\cite{atkey2009} encode protocol states at
the type level but require non-standard syntax and are confined to a single
domain.

\paragraph{The gap.}
Future conditions themselves are not new.
\citet{song2025futurecond} introduce them as a first-class extension to
pre/post contracts and give a logical semantics;
\citet{song2024provenfix} encode temporal properties as regular expressions
and use derivatives to drive automated program repair in the ProveNFix tool.
Neither, however, treats an obligation as a \emph{value} that a program can
carry, combine, and hand on.
The first works at the level of a logical system, with no implementation;
the second embeds its reasoning inside a whole-program analysis that a
programmer cannot invoke compositionally, and that is specific to traces.
What is missing is an account of how an obligation propagates through
effectful code one bind at a time --- and, with it, a library a Haskell
programmer can actually use.

\paragraph{Why a library, and not a type system or a model checker.}
Temporal properties have been attacked from many directions: encoded in
indexed or graded types, discharged by a model checker, or proved in a
theorem prover.
Each buys stronger guarantees at a cost.
Type-level encodings force the protocol into the type of every
intermediate computation, so a change to the specification perturbs
type signatures throughout the program.
Model checkers and provers work on a model of the program rather than the
program, and demand that the property be fixed before the code is written.
Making the obligation an ordinary Haskell \emph{value} costs the static
guarantee, but buys three things those approaches do not offer together:
obligations can be computed --- built from runtime data, passed to
functions, stored in data structures;
one bind formula serves every domain, so a new obligation language is a new
\lstinline{Composable} instance rather than a new encoding;
and the unmet remainder is itself an inspectable value, so a failure
\emph{names} what is still owed instead of reporting that a proof did not
go through.

\paragraph{What Pledge does, and what it does not.}
A \lstinline{Pledge m eff a} value is a computation in the underlying monad
\lstinline{m} that additionally produces three annotation fields
$(\pre, \post, \fut) :: \mathit{eff}$, and composing such values propagates
those annotations algebraically alongside whatever \lstinline{m} does.
Two consequences are worth stating before the details.

First, the annotations are \emph{path-sensitive}.
Because \lstinline{m} is a monad, bind must run the first action before it
can apply the continuation to the result, so the annotations that come out
describe the one path the program actually took --- not an
over-approximation of every path it might have taken.
This is what lets \S\ref{sec:ltl} read \lstinline{Pledge} as a
\emph{runtime monitor}: the residual after each bind is the monitor's state,
computed beside a real effect handler rather than against a static model.

Second, nothing intervenes, as the library stands today.
The residual is a value to be inspected; a program that fails to discharge
an obligation still runs to completion, and \lstinline{Pledge} names the
omission rather than preventing it.
Promoting the residual to a runtime guard is future work we scope precisely
in \S\ref{sec:usability}.

\paragraph{Our approach.}
A value of type \lstinline{Pledge m eff a} carries three annotations:
a precondition $\pre :: \mathit{eff}$ (what must have held before),
a postcondition $\post :: \mathit{eff}$ (what this action produces), and
a future condition $\fut :: \mathit{eff}$ (what must hold afterwards).
Monadic bind propagates these through two quotient operations on
\lstinline{eff}:
\begin{align}
  \pre(p \mathbin{>\!\!>\!\!=} q)  &\;=\; \pre(p) \sqcap (\pre(q) \rquot \post(p))
  \label{eq:pre-bind}  \\
  \post(p \mathbin{>\!\!>\!\!=} q) &\;=\; \post(p) \cdot \post(q)
  \label{eq:post-bind} \\
  \fut(p \mathbin{>\!\!>\!\!=} q)  &\;=\; (\fut(p) \lquot \post(q)) \sqcap \fut(q)
  \label{eq:fut-bind}
\end{align}
A fully-discharged residual --- $\pre = \univ$ and $\fut = \univ$ --- is the
certificate that the composed program is contract-correct;
anything else names what is still owed.
\S\ref{sec:pledge} develops this formula and the six-operation
\lstinline{Composable} algebra it is stated over; \S\ref{sec:monad-laws}
identifies exactly what that algebra must satisfy for
Equations~\eqref{eq:pre-bind}--\eqref{eq:fut-bind} to obey the monad laws;
\S\ref{sec:re-instance} instantiates the algebra with regular expressions
and Brzozowski derivatives and proves the instance sound;
\S\ref{sec:ltl}--\S\ref{sec:mtl} build two front-ends on top of that
instance, and \S\ref{sec:gre} builds a fourth, SMT-checked instance beside
it; \S\ref{sec:usability}--\S\ref{sec:experiments} lay out, concretely, the
three pieces of engineering the design still needs.

\paragraph{Contributions.}
\begin{enumerate}
  \item \textbf{The \lstinline{Pledge} monad transformer and the
        \lstinline{Composable} algebra} (\S\ref{sec:pledge}): six operations
        that make future conditions first-class Haskell values propagated
        by a single bind formula, closed under products and pointwise
        lifting through any monad.

  \item \textbf{Mechanised metatheory, and its exact boundary}
        (\S\ref{sec:monad-laws}--\S\ref{sec:re-instance}).
        We identify twelve algebraic axioms sufficient for the three monad
        laws on the $(\ret, \post, \fut)$ triple and prove the laws
        \lstinline{sorry}-free in Lean~4.
        Mechanising the $\RE$ instance against those axioms then
        \emph{delimits} them, and we state and prove the two side
        conditions --- single-word postconditions, past-closed
        preconditions --- that restore full equality, together with a
        soundness theorem: a fully discharged $\RE$ residual entails that
        the emitted trace met every obligation its primitives named.

  \item \textbf{Two front-ends on top of $\RE$}
        (\S\ref{sec:ltl}, \S\ref{sec:mtl}): an $\LTL$ layer that reframes
        \lstinline{Pledge} as a compositional runtime monitor, and a
        bounded $\MTL$ layer adding discrete-time intervals, both compiling
        to $\RE$ so \S\ref{sec:re-instance}'s theory applies unchanged.

  \item \textbf{A fourth, SMT-checked instance} (\S\ref{sec:gre}): $\GRE$
        pairs an $\RE$ with a Presburger predicate over heap addresses,
        discharged by Z3, and because that predicate is checked
        symbolically it gives \lstinline{Pledge} a static-verification mode
        distinct from $\LTL$'s runtime one.

  \item \textbf{A concrete roadmap for the unfinished third of the design}
        (\S\ref{sec:usability}--\S\ref{sec:experiments}): making a
        violation's location traceable, automating the Lean proof
        obligations a new instance incurs, and the empirical evaluation the
        design has not yet had.
\end{enumerate}

%% =========================================================================
\section{What Is Pledge}
\label{sec:pledge}
%% =========================================================================

\subsection{The \texttt{Composable} Class}
\label{sec:composable}

We abstract the specification algebra via a type class with six operations:
\begin{lstlisting}
class Composable a where
  concatenation :: a -> a -> a   -- (·), identity: empty
  conjunction   :: a -> a -> a   -- (⊓), identity: universe
  empty         :: a             -- identity for (·)
  universe      :: a             -- identity for (⊓)
  leftQuotient  :: a -> a -> a   -- dividend ∖ divisor
  rightQuotient :: a -> a -> a   -- dividend ∕ divisor
\end{lstlisting}
The library provides Unicode infix aliases:
\lstinline{(·) = concatenation},
\lstinline{(⊓) = conjunction},
\lstinline{a ∖ b = leftQuotient b a}, and
\lstinline{p ∕ q = rightQuotient q p}.
The swapped arguments in the infix definitions are intentional:
$\fut \lquot \post$ means ``left-quotient of $\fut$ by $\post$,'' i.e.,
strip $\post$ from the left of $\fut$;
$\pre \rquot \post$ means ``right-quotient of $\pre$ by $\post$,''
i.e., strip $\post$ from the right of $\pre$.

\paragraph{Algebraic roles.}
\lstinline{concatenation} sequences effects;
\lstinline{conjunction} combines simultaneous constraints;
\lstinline{empty} and \lstinline{universe} appear in \lstinline{pure}.
The two quotient operations implement \emph{obligation discharge}:
$\fut \lquot \post$ asks what of the original future obligation remains
unmet once the continuation has posted $\post$, and $\pre \rquot \post$
asks what portion of the next action's precondition is not yet satisfied
by the current action's postcondition.
Two base cases anchor both quotients:
$r \lquot \varepsilon = r$ (nothing observed, obligation unchanged) and
$r \lquot \top = \top$ (universal postcondition covers anything).
The symmetric laws hold for $\rquot$.

\paragraph{Composing instances.}
The instance for pairs lifts \lstinline{Composable} component-wise, and a
separate instance lifts it pointwise through any \lstinline{Applicative}
functor \lstinline{m}, so \lstinline{Composable (m eff)} holds whenever
\lstinline{Composable eff} does:
\begin{lstlisting}
instance (Composable a, Composable b) => Composable (a, b) where
  concatenation (a1,b1) (a2,b2) = (concatenation a1 a2, concatenation b1 b2)
  leftQuotient  (a1,b1) (a2,b2) = (leftQuotient  a1 a2, leftQuotient  b1 b2)
  -- etc.
\end{lstlisting}
Independent specification dimensions --- a trace obligation and a heap
obligation, say --- combine by pairing their \lstinline{Composable}
instances; \S\ref{sec:gre} instead gives a fourth instance that couples the
two dimensions directly.

\subsection{The \texttt{Pledge} Monad Transformer}
\label{sec:type}

\begin{lstlisting}
newtype Pledge m eff a = Pledge { runPledge :: m (a, eff, eff, eff) }
\end{lstlisting}
A value of type \lstinline{Pledge m eff a} is a computation in the
underlying monad \lstinline{m} that produces a 4-tuple
$(a, \pre, \post, \fut)$: the return value and three annotation fields.
When \lstinline{m = IO}, the annotations can be built from actual runtime
values (e.g., a freshly allocated file handle), because the 4-tuple is
produced by the \lstinline{IO} action itself.

Two combinators support interoperability:
\begin{lstlisting}
-- lift a plain m-action with trivial annotations
liftPledge :: (Composable eff, Applicative m) => m a -> Pledge m eff a
liftPledge ma = Pledge $ fmap (, universe, empty, universe) ma

-- run once and collect all four components
inspect :: Functor m => Pledge m eff a -> m (PledgeResult eff a)
inspect (Pledge ma) = fmap (\(a,pre,post,fut) -> PledgeResult a pre post fut) ma
\end{lstlisting}
\lstinline{liftPledge} embeds any \lstinline{m}-action as a \lstinline{Pledge}
with no precondition, no postcondition, and no future obligation.
\lstinline{inspect} runs the action exactly once and packages the result;
callers should prefer \lstinline{inspect} over the individual field accessors
(\lstinline{getPre}, \lstinline{getPost}, \lstinline{getFut}) when
\lstinline{m} has side effects.

\subsection{Monad Instance}

\lstinline{pure} introduces a value with no effect and no obligation, and
bind sequences two computations and propagates annotations according to
Equations~\eqref{eq:pre-bind}--\eqref{eq:fut-bind}:
\begin{lstlisting}
pure x = Pledge $ pure (x, universe, empty, universe)

Pledge ma >>= g = Pledge $ do
    (a, preA, postA, futA) <- ma
    (b, preB, postB, futB) <- runPledge (g a)
    return (b,  preA ⊓ (preB ∕ postA),
                postA · postB,
               (futA ∖ postB) ⊓ futB)
\end{lstlisting}
When the future component normalises to $\top$, all future obligations are
discharged.
When $\post(A)$ does not cover $\pre(B)$, the residual precondition is
$\varnothing$, flagging the violation; when it fully covers $\pre(B)$, the
residual is $\varepsilon$ and the combined precondition reduces to
$\pre(A)$, matching the standard Hoare sequencing rule.
\S\ref{sec:monad-laws} shows that, despite this asymmetric phrasing, all
four fields are treated by exactly the same algebra, and identifies the
condition under which the triple $(\ret, \post, \fut)$ satisfies the monad
laws with equality.

\paragraph{One operation, two directions.}
The two quotients are not independent machinery.
A future condition is an obligation on the trace still to come; a
precondition is an obligation on the trace already past; and reversing time
exchanges them.
Concretely, in the $\RE$ instance of \S\ref{sec:re-instance} the right
quotient \emph{is} the left quotient conjugated by reversal,
\[
  r \rquot s \;=\; \mathrm{rev}\bigl(\mathrm{rev}(r) \lquot \mathrm{rev}(s)\bigr),
\]
the same regular expression $\Sigma^* \cdot e \cdot \Sigma^*$ serves as
``$e$ happened earlier'' in a $\pre$ slot and ``$e$ must happen later'' in a
$\fut$ slot, and the four right-quotient axioms of \S\ref{sec:monad-laws}
are the four left-quotient axioms with the direction of composition flipped.
This symmetry is what makes $\pre$ cheap --- it reuses the machinery
$\fut$ already needs --- and it also explains the one place the analogy
stops: monadic bind is \emph{oriented}, composing forward, so the
backward-facing field cannot satisfy the monad laws in the same
orientation.
We return to this in \S\ref{sec:monad-laws}.

\subsection{How the Specification Meets Real Effects}
\label{sec:effects}

Because annotations live inside the underlying \lstinline{m}, the
\lstinline{Pledge} monad transformer integrates naturally with real
effectful code, in one of two ways.
When the specification primitives use
\lstinline{Pledge \$ return (val, pre, post, fut)}, annotations are pure
Haskell values and the underlying action does nothing; composing such
primitives with \lstinline{>>=} builds a specification term the programmer
inspects independently of running any real effect.
When the underlying action is real, the annotations can instead reference
actual runtime values:
\begin{lstlisting}
openFile :: FilePath -> IO.IOMode -> Pledge IO (RE IO.Handle) IO.Handle
openFile fn mode = Pledge $ do
    h <- IO.openFile fn mode            -- actual IO
    return (h, universe,
            Single (Atom "open" h),     -- post uses the real handle h
            finally (Atom "close" h))   -- fut requires closing h specifically
\end{lstlisting}
Here \lstinline{h :: IO.Handle} is the actual file handle produced at
runtime, embedded directly into the \lstinline{RE IO.Handle} annotations.
The left-quotient machinery then propagates this concrete handle through
subsequent binds, so a residual \lstinline{finally(close h)} names the
specific unclosed handle --- the mechanism that lets \S\ref{sec:ltl} treat
\lstinline{inspect} as running a monitor over a real execution rather than
a specification detached from it.

%% =========================================================================
\section{What Makes Pledge Follow the Monad Law}
\label{sec:monad-laws}
%% =========================================================================

Every monad must satisfy three equational laws.
This section states, abstractly over the \lstinline{Composable} algebra,
exactly what suffices for \lstinline{Pledge} to do so --- and is equally
precise about what that statement excludes.
The exclusion is structural and expected: the $\pre$ field is oriented
against the direction of bind, and is carried beside the lawful triple
rather than inside it.
\S\ref{sec:re-instance} then asks, for the concrete $\RE$ instance,
\emph{which} of the axioms below actually hold, and under what condition on
how annotations are written.

%% -------------------------------------------------------------------------
\subsection{Why \texorpdfstring{$\pre$}{pre} sits outside the triple}
\label{sec:orientation}

Recall from \S\ref{sec:pledge} that $\pre$ and $\fut$ are the same notion
pointing in opposite directions, and $\rquot$ is $\lquot$ conjugated by
reversal.
The monad laws break that symmetry, because bind does not respect it:
$p \mathbin{>\!\!>\!\!=} q$ runs $p$ and \emph{then} $q$, so every law
equates rearrangements of a forward-composed sequence.
A field that accumulates forwards can satisfy such laws; a field that
accumulates backwards cannot be expected to, and $\pre$ does not.

Concretely, left identity requires the annotations of
$\mathtt{pure}\;x \mathbin{>\!\!>\!\!=} f$ to equal those of $f\,x$.
For $\fut$ this holds because $\mathtt{pure}$ contributes $\top$ on the
left, and $\top$ is the unit that $\lquot$ preserves (axiom L2 below).
For $\pre$ the corresponding step needs $\top \rquot \post(f\,x) = \top$,
which is axiom R2 --- the \emph{mirror} of L2, and true for exactly the
same reason.
So $\pre$ is not lawless; it satisfies the mirror-image laws.
What it cannot do is satisfy the forward-oriented laws simultaneously with
$\fut$ under a single equality, because a detected precondition violation
must stay visible under further binding rather than being normalised away
--- a Hoare-rule discipline, not an oversight.

We therefore prove the three laws for the triple $(\ret, \post, \fut)$ and
carry $\pre$ alongside, in the manner of a writer-monad log: inspected,
propagated by Equation~\eqref{eq:pre-bind}, and not required to satisfy the
monad laws on its own.
All four fields \emph{are} covered by the axioms below; the asymmetry is in
which equations we then ask them to satisfy, not in how they are treated
algebraically.

%% -------------------------------------------------------------------------
\subsection{Twelve axioms}
\label{sec:axioms}

The laws hold assuming the twelve axioms in Table~\ref{tab:axioms},
stated abstractly over the \lstinline{Composable} algebra.
The reversal duality is visible in the table: R1--R4 are L1--L4 with the
order of composition reversed, and nothing else.

\begin{table}[ht]
\caption{The twelve \lstinline{Composable} axioms required to prove the
         monad laws for the $(\ret, \post, \fut)$ triple.}
\label{tab:axioms}
\small
\begin{tabular}{lll}
\toprule
Label & Axiom & Used for \\
\midrule
S1 & $\varepsilon \cdot a = a$ & Law 1, post \\
S2 & $a \cdot \varepsilon = a$ & Law 2, post \\
S3 & $(a \cdot b) \cdot c = a \cdot (b \cdot c)$ & Law 3, post \\
\midrule
C1 & $a \sqcap b = b \sqcap a$ & Law 2, pre / fut \\
C2 & $(a \sqcap b) \sqcap c = a \sqcap (b \sqcap c)$ & Law 3, pre / fut \\
C3 & $\top \sqcap a = a$ & Laws 1, 2 \\
\midrule
L1 & $r \lquot \varepsilon = r$ & Law 2, fut \\
L2 & $\top \lquot r = \top$ & Law 1, fut \\
L3 & $r \lquot (a \cdot b) = (r \lquot a) \lquot b$ & Law 3, fut \\
L4 & $(a \sqcap b) \lquot c = (a \lquot c) \sqcap (b \lquot c)$ & Law 3, fut \\
\midrule
R1 & $a \rquot \varepsilon = a$ & Law 1, pre \\
R2 & $\top \rquot r = \top$ & Law 2, pre \\
R3 & $(a \rquot b) \rquot c = a \rquot (c \cdot b)$ & Law 3, pre \\
R4 & $(a \sqcap b) \rquot c = (a \rquot c) \sqcap (b \rquot c)$ & Law 3, pre \\
\bottomrule
\end{tabular}
\end{table}

\subsection{The three laws}
\label{sec:three-laws}

We verify all three laws using these axioms.
In each case we show that both sides of the law produce equal
$(\ret, \post, \fut)$ triples.
It is worth noting which axioms each law consumes, because
\S\ref{sec:boundary} shows that not all twelve are unconditionally
available for a given instance:
left identity uses S1, C3, L2, R1;
right identity uses S2, C1, C3, L1, R2;
and associativity uses S3, C2, L3, L4, R3, R4.

\paragraph{Left identity: \lstinline{pure x >>= f = f x}.}
Let $p = \mathtt{pure}\;x = (x, \top, \varepsilon, \top)$ and $q = f\,x$.
By Equations~\eqref{eq:pre-bind}--\eqref{eq:fut-bind}:
\begin{align*}
  \ret  &: \ret(q). \checkmark \\
  \post &: \varepsilon \cdot \post(q) = \post(q). \quad\text{(S1)} \checkmark \\
  \fut  &: (\top \lquot \post(q)) \sqcap \fut(q)
         = \top \sqcap \fut(q) = \fut(q). \quad\text{(L2, C3)} \checkmark
\end{align*}
For $\pre$: $\pre(q) \rquot \varepsilon = \pre(q)$ (R1), so $\pre = \top \sqcap \pre(q)
= \pre(q)$ (C3). $\checkmark$

\paragraph{Right identity: \lstinline{e >>= pure = e}.}
Let $q = \mathtt{pure}(\ret(e)) = (\ret(e), \top, \varepsilon, \top)$.
\begin{align*}
  \ret  &: \ret(e). \checkmark \\
  \post &: \post(e) \cdot \varepsilon = \post(e). \quad\text{(S2)} \checkmark \\
  \fut  &: (\fut(e) \lquot \varepsilon) \sqcap \top
         = \fut(e) \sqcap \top = \top \sqcap \fut(e) = \fut(e). \quad\text{(L1, C1, C3)} \checkmark
\end{align*}

\paragraph{Associativity.}
Let $r_1 = \ret(e)$, $r_2 = \ret(f\,r_1)$.
Both sides yield $\ret(g\,r_2)$ and
$\post(e) \cdot \post(f\,r_1) \cdot \post(g\,r_2)$ (by S3).
For the future, unfolding Equation~\eqref{eq:fut-bind} on the LHS and
applying L4 then L3 gives:
\[
  \bigl((\fut(e) \lquot \post(f\,r_1)) \lquot \post(g\,r_2)\bigr) \sqcap
  (\fut(f\,r_1) \lquot \post(g\,r_2)) \sqcap \fut(g\,r_2)
\]
which equals the RHS by C2 and L3. $\checkmark$
For $\pre$, applying R4 then R3 gives the same result on both sides by C2. $\checkmark$

The three proofs above are mechanised in Lean~4 in
\texttt{Formalization/PledgeMonadLaws.lean}, using an abstract
\lstinline{Composable} structure parametrised over the twelve axioms of
Table~\ref{tab:axioms}.
The file carries no \lstinline{sorry}; all proofs are term-mode derivations
from the stated axioms, and depend on Mathlib~\cite{mathlib4} only for
basic list and set lemmas.
Because the theorem is proved once, abstractly, \emph{any} instance that
discharges the twelve axioms inherits the monad laws for free ---
\S\ref{sec:re-instance} is where we pay that debt for $\RE$.

%% =========================================================================
\section{RE Satisfies the Monad Laws}
\label{sec:re-instance}
%% =========================================================================

\S\ref{sec:monad-laws} proved the monad laws \emph{relative to} twelve
axioms.
This section discharges those axioms for a concrete instance: extended
regular expressions over a typed event alphabet, with Brzozowski
derivatives computing both quotients.
Mechanising this instance in Lean does not merely confirm the axioms ---
it \emph{delimits} them, and the delimitation is the main technical content
of this section.

\subsection{Extended Regular Expressions over Typed Events}
\label{sec:re}

Events are the atomic units of observation.
An \lstinline{Event t} carries a name and a typed payload, or is a
wildcard, or records a bare \emph{use} of a value without naming an
operation (for obligations that are about touching a resource at all,
rather than any one named event on it):
\begin{lstlisting}
data Event t = Atom String t   -- concrete event, e.g. Atom "acquire" (Num 1)
             | Usage t          -- any access to a value, unnamed
             | Wildcard         -- matches any single event
\end{lstlisting}
The type parameter \lstinline{t} allows the payload to be any type,
including runtime values such as file handles or heap addresses; the
concrete payload type used throughout the examples below, \lstinline{Values},
also underlies the Presburger terms of \S\ref{sec:gre}.
The specification language is an extended regular expression type:
\begin{lstlisting}
data RE t = Bot | Epsilon | Single (Event t)
          | Seq (RE t) (RE t) | Or (RE t) (RE t)
          | And (RE t) (RE t) | Star (RE t) | Not (RE t)
\end{lstlisting}
The constructors are standard: $\varnothing$ (\lstinline{Bot}),
$\varepsilon$ (\lstinline{Epsilon}), single-event match
(\lstinline{Single}), concatenation (\lstinline{Seq}), union
(\lstinline{Or}), intersection (\lstinline{And}), Kleene star
(\lstinline{Star}), and complement (\lstinline{Not}).
\lstinline{And} is a primitive constructor so its derivative is computed
directly rather than through three extra \lstinline{Not} nodes.

Four derived forms appear throughout the paper:
\begin{align*}
  \anything &\;\triangleq\; \neg\varnothing
            &&\text{(universal language, written $\top$)} \\
  \finally{e} &\;\triangleq\; \anything \cdot e \cdot \anything
            &&\text{($e$ must eventually occur)} \\
  \never{e} &\;\triangleq\; \neg\finally{e}
            &&\text{($e$ must never occur)} \\
  \noUntil{e}{g} &\;\triangleq\;
            \neg\bigl((\Sigma\setminus\{g\})^* \cdot e \cdot \anything\bigr)
            &&\text{($e$ must not occur before $g$)}
\end{align*}
\noindent
$\previously{e}$ is an alias for $\finally{e}$ used in \lstinline{pre} slots
to assert that $e$ occurred somewhere in the preceding trace; the underlying
RE is identical, but the name signals intent, and \S\ref{sec:pre-condition}
is precisely about why this alias matters.

\subsection{Algebraic Complement and Brzozowski Derivatives}
\label{sec:deriv}

The Brzozowski derivative~\cite{brzozowski1964} $\partial_a(r)$ computes
the language of continuations after reading event $a$.
The key law for complement is:
\begin{equation}
  \partial_a(\neg r) \;=\; \neg\bigl(\partial_a(r)\bigr)
  \label{eq:deriv-not}
\end{equation}
Together with the De Morgan laws and double-negation elimination, this lets
complement be propagated purely algebraically, without building a DFA.
The nullability function $\nu(r)$ ($= \mathtt{True}$ iff $\varepsilon \in L(r)$)
satisfies $\nu(\neg r) = \neg\,\nu(r)$:
\begin{lstlisting}
nullable :: RE t -> Bool
nullable (Not r)     = not (nullable r)
nullable (Star _)    = True
nullable (Seq r1 r2) = nullable r1 && nullable r2
nullable (Or  r1 r2) = nullable r1 || nullable r2
nullable (And r1 r2) = nullable r1 && nullable r2
nullable Epsilon     = True
nullable _           = False

derivative :: Eq t => Event t -> RE t -> RE t
derivative e (Not r)      = Not (derivative e r)   -- da(Not r) = Not(da(r))
derivative e (Single p)   = if subsumesEvent e p then Epsilon else Bot
derivative e (Seq r1 r2)
    | nullable r1 = Or (Seq (derivative e r1) r2) (derivative e r2)
    | otherwise   = Seq (derivative e r1) r2
derivative e (Or  r1 r2)  = Or  (derivative e r1) (derivative e r2)
derivative e (And r1 r2)  = And (derivative e r1) (derivative e r2)
derivative e (Star r)     = Seq (derivative e r) (Star r)
derivative _ _            = Bot
\end{lstlisting}
Complement requires care when computing which events can begin a word:
naively returning $\mathit{first}(\neg r, \Sigma_0) = \mathit{first}(r, \Sigma_0)$
is unsound.
The correct characterisation, over a finite working alphabet $\Sigma_0$, is
\begin{equation}
  e \in \mathit{first}(\neg r,\,\Sigma_0)
  \;\iff\;
  e \in \Sigma_0
  \;\wedge\;
  \norm{\partial_e(r)} \neq \univ,
  \label{eq:first-not}
\end{equation}
that is, $e$ can open a word in $L(\neg r)$ iff, after consuming $e$,
the residual $\partial_e(r)$ does not cover the entire language.
The working alphabet itself must include \lstinline{Wildcard} as the
representative of ``some event other than the named ones''; omitting it is
not a conservative simplification but a soundness bug, since a pair of
operands naming no concrete event at all (such as $\Sigma^*$) would then
have an empty working alphabet and quotient to $\varnothing$ regardless of
its actual denotation.

\subsection{The \texttt{Composable RE} Instance}
\label{sec:composable-re}

\begin{lstlisting}
instance Eq t => Composable (RE t) where
  concatenation = Seq
  conjunction   = And
  empty         = Epsilon
  universe      = Not Bot        -- = Sigma*
  leftQuotient  = reLeftQuotient
  rightQuotient = reRightQuotient
\end{lstlisting}
The left-quotient $r_2 \lquot r_1$ is computed by a worklist traversal over
reachable pairs of Brzozowski derivatives, which keeps termination
tractable and is what \lstinline{reLeftQuotient} implements.
The right-quotient is computed via reversal,
$r_2 \rquot r_1 = \mathit{rev}(\mathit{rev}(r_2) \lquot \mathit{rev}(r_1))$,
where $\mathit{rev}(r)$ reverses the sequences accepted by $r$ --- the
mechanical form of the reversal duality of \S\ref{sec:pledge}.
A normaliser applies the laws in Table~\ref{tab:normalize-laws} to simplify
RE terms, including $\Sigma^* \cdot r = \Sigma^*$ for nullable $r$, without
which a fully-discharged residual such as $\Sigma^* \cdot (a \cdot \Sigma^*
\vee \varepsilon)$ would stay syntactically distinct from $\top$ and be
reported as outstanding.

\begin{table}[ht]
\caption{Normalisation laws applied by \lstinline{normalize}, recursively
         bottom-up.}
\label{tab:normalize-laws}
\begin{tabular}{ll}
\toprule
Law & Constructor \\
\midrule
$\varepsilon \cdot r = r$, $r \cdot \varepsilon = r$, $\top\!\cdot\!r = \top$ (r nullable) & \lstinline{Seq} \\
$\varnothing \vee r = r$, $r \vee \top = \top$ & \lstinline{Or} \\
$\varnothing \wedge r = \varnothing$, $r \wedge \top = r$ & \lstinline{And} \\
$\neg\neg r = r$, De Morgan, $\neg\varnothing = \top$, $\neg\top = \varnothing$ & \lstinline{Not} \\
$\varnothing^* = \varepsilon$, $\varepsilon^* = \varepsilon$ & \lstinline{Star} \\
\bottomrule
\end{tabular}
\end{table}

\paragraph{Termination.}
The quotient recursion terminates because the set of distinct normalised
derivatives of any regular expression is finite modulo associativity,
commutativity and idempotence of \lstinline{Or}/\lstinline{And} --- a
classical result of Brzozowski~\cite{brzozowski1964}, extended to
complement by Owens et al.~\cite{owens2009}.
Cycle detection in the worklist must therefore compare states modulo those
identities, not structurally: a divisor whose derivative does not shrink
under structural comparison, such as $\Sigma^*$, would otherwise make the
recursion diverge.

%% -------------------------------------------------------------------------
\subsection{What the \texorpdfstring{$\RE$}{RE} instance actually satisfies}
\label{sec:boundary}

Eight of the twelve axioms hold outright for $\RE$; four do not
unconditionally, and Table~\ref{tab:axiom-status} records which.

\begin{table}[ht]
\caption{Status of the twelve axioms for the $\RE$ instance, as mechanised
         in \texttt{Formalization/REComposableLaws.lean}.}
\label{tab:axiom-status}
\small
\begin{tabular}{lll}
\toprule
Axioms & Status & Condition \\
\midrule
S1--S3, C1--C3 & hold & --- \\
L1, L3, R1, R3 & hold & --- \\
L2, R2         & hold conditionally & $L(\text{divisor}) \neq \varnothing$ \\
L4, R4         & one inclusion only & $\subseteq$ holds; $\supseteq$ fails \\
\bottomrule
\end{tabular}
\end{table}

\paragraph{Why L2 and R2 need a non-empty divisor.}
$\top \lquot r$ collects every suffix left over after consuming some word of
$L(r)$.
If $L(r) = \varnothing$ there is no such word and the quotient is
$\varnothing$, not $\top$: a postcondition denoting the empty language would
assert that the action produced an impossible trace.

\paragraph{Why L4 and R4 hold in only one direction.}
Written out, L4 asks
$(a \sqcap b) \lquot c = (a \lquot c) \sqcap (b \lquot c)$.
The left-hand side demands \emph{one} word $u \in L(c)$ that works for both
conjuncts; the right-hand side lets each conjunct choose its own.
Taking $L(c) = \{a, b\}$, $L(a) = \{a\}$ and $L(b) = \{b\}$: on the left,
$L(a) \cap L(b) = \varnothing$, so the quotient is $\varnothing$; on the
right, the first conjunct picks $u = a$ and the second picks $u = b$, each
leaving $\varepsilon$, so the meet is $\{\varepsilon\}$.
Only $\subseteq$ survives.

\paragraph{Consequence: associativity is an inclusion.}
Associativity consumes L4 and R4, so in general it holds only as
\[
  \fut\bigl((p \mathbin{>\!\!>\!\!=} f) \mathbin{>\!\!>\!\!=} g\bigr)
  \;\subseteq\;
  \fut\bigl(p \mathbin{>\!\!>\!\!=} (\lambda x.\, f\,x \mathbin{>\!\!>\!\!=} g)\bigr),
\]
and symmetrically for $\pre$, the left-nested bracketing yielding the
\emph{stronger} obligation.
This is not a vacuous gap: over a pool of nine representative annotations in
each of five positions, the two bracketings differ in $902$ of $59{,}049$
combinations, strictly in the direction stated every time.

\subsection{Single-word postconditions restore equality}
\label{sec:single-word}

The gap closes under a condition that costs nothing in practice.

\begin{lemma}[Unique witness]
\label{lem:unique-witness}
If $L(c)$ is a single word $\{u\}$, then L4 and R4 hold with equality, and
L2 and R2 hold.
\end{lemma}
\begin{proof}
$\{u\}$ is non-empty, giving L2 and R2.
For L4, $(a \sqcap b) \lquot \{u\} = \{w \mid uw \in L(a) \cap L(b)\}
= \{w \mid uw \in L(a)\} \cap \{w \mid uw \in L(b)\}
= (a \lquot \{u\}) \sqcap (b \lquot \{u\})$;
the choice of witness is forced, so the two sides cannot diverge.
R4 is the mirror image.
\end{proof}

\begin{theorem}[Laws for single-word postconditions]
\label{thm:single-word}
If every primitive's $\post$ denotes a single word, then all twelve axioms
hold for the $\RE$ instance and the three monad laws hold with equality.
\end{theorem}
\begin{proof}
Every $\lquot$ and $\rquot$ in Equations~\eqref{eq:pre-bind}
and~\eqref{eq:fut-bind} has a $\post$ field as its divisor, so
Lemma~\ref{lem:unique-witness} applies at every use; the remaining eight
axioms hold unconditionally.
\end{proof}

The hypothesis is not a restriction, because it is what the library
produces anyway and is preserved by composition: a primitive emits exactly
one event; \lstinline{pure} and \lstinline{liftPledge} emit $\varepsilon$;
and by Equation~\eqref{eq:post-bind} the postcondition of a composite is
$\post(p) \cdot \post(q)$, a concatenation of two single-word languages and
hence itself a single word.
What the condition rules out is a primitive that advertises a \emph{choice}
of postconditions --- \lstinline{Or}, \lstinline{Star} or a complement in a
$\post$ slot; such an annotation is expressible and silently costs
associativity, which \S\ref{sec:usability} returns to as an API-hardening
opportunity.

\subsection{Preconditions must describe the whole past}
\label{sec:pre-condition}

A second side condition governs $\pre$, and violating it is the more
insidious failure of the two, because it produces a confident report of a
contract violation for a correct program.

The natural way to annotate \lstinline{release} is ``an \lstinline{acquire}
must have happened'', and the natural way to write that is
$\pre = \{\mathtt{acquire}\}$ --- the singleton language.
That is wrong, and not for the reason one expects: it says the acquire must
be the \emph{immediately preceding} event, and such a condition does not
survive composition.
Unfolding Equation~\eqref{eq:pre-bind} for
$p \mathbin{>\!\!>\!\!=} q$ with $\pre(q) = \{e\}$ and $\post(p) = \{e\}$:
\[
  \pre(q) \rquot \post(p) \;=\; \{w \mid w \cdot e \in \{e\}\}
                          \;=\; \{\varepsilon\},
\]
so $\pre(p \mathbin{>\!\!>\!\!=} q) = \pre(p) \sqcap \{\varepsilon\}$,
which is $\varnothing$ whenever $\pre(p)$ does not accept the empty trace.
The obligation was discharged, and the residual nevertheless reports a
violation.

The fix is to require preconditions to be closed under extension on the
left --- properties of the entire preceding trace, not of its last event.

\begin{definition}[Past-closed]
A language $P$ is \emph{past-closed} if $P = \Sigma^* \cdot P$: whenever
$w \in P$ and $u \in \Sigma^*$, also $uw \in P$.
\end{definition}

\begin{proposition}[Closure]
\label{prop:past-closed}
Past-closed languages contain $\top$ and are closed under $\sqcap$ and
under $\rquot$ by an arbitrary divisor.
Consequently, if every primitive's $\pre$ is past-closed, so is the $\pre$
of every program composed from them.
\end{proposition}
\begin{proof}
$\top = \Sigma^*$ is past-closed.
For $\sqcap$: $\Sigma^*(A \cap B) \subseteq \Sigma^*A \cap \Sigma^*B
= A \cap B$, and the reverse inclusion takes $u = \varepsilon$.
For $\rquot$: let $w \in P \rquot s$, so $wu \in P$ for some $u \in L(s)$;
for any $v$, past-closure of $P$ gives $vwu \in P$, hence
$vw \in P \rquot s$.
The final claim follows by induction on program structure, since
Equation~\eqref{eq:pre-bind} builds $\pre$ from $\sqcap$ and $\rquot$ alone.
\end{proof}

The $\previously{e} = \Sigma^* \cdot e \cdot \Sigma^*$ combinator of
\S\ref{sec:re} is past-closed, and is the intended way to write a
precondition; the singleton $\{e\}$ is not.
With it, the collapse above does not occur: $\previously{e} \rquot \{e\}
= \Sigma^*$, the obligation is discharged to $\top$, and the enclosing
$\sqcap$ leaves $\pre(p)$ untouched.
Proposition~\ref{prop:past-closed} guarantees this is stable under further
composition.

Neither this condition nor Theorem~\ref{thm:single-word} is enforced by the
type \lstinline{eff}, and both are easy to violate by writing the
annotation that first comes to mind; \S\ref{sec:usability} discusses making
them unwriteable rather than merely documented.

\subsection{Mechanisation and validation}
\label{sec:mechanisation}

A companion Lean file, \texttt{Formalization/REComposableLaws.lean},
develops the $\RE$ instance at the level of languages: it proves the eight
unconditional axioms, proves L2 and R2 under the non-emptiness hypothesis,
proves the $\subseteq$ direction of L4 and R4 and exhibits the
counterexample to the converse, and proves derivative and
\lstinline{nullable} correctness.
It carries no \lstinline{sorry}.

We complement the mechanisation with exhaustive testing over bounded pools,
which is what first exposed the boundary.
Quotients are checked against their set-theoretic definition
$w \in r \lquot s \iff \exists u \in L(s).\, uw \in L(r)$
over all pairs from a pool of $90$ regular expressions and all words up to
length two --- $56{,}700$ checks for each quotient, with no discrepancy.
The axioms are then checked over the same pool, reproducing
Table~\ref{tab:axiom-status} exactly and independently of the Lean
development.
Finally, restricting postconditions to single words as
Theorem~\ref{thm:single-word} requires, and letting futures and
preconditions range over $\RE$ terms including complement, star and
intersection, associativity holds in all $43{,}200$ combinations tested ---
against the $902$ failures found without the restriction.

\subsection{Soundness}
\label{sec:soundness}

A fully-discharged residual ($\univ$ in $\fut$, $\univ$ in $\pre$) serves as
a certificate of contract-correctness.
We now prove that this certificate is sound.
For an RE $r$ and a finite trace $w$, write $w \models r$ for $w \in L(r)$
and $\partial_w(r)$ for the iterated derivative.

\begin{lemma}[Derivative correctness]
\label{lem:deriv-correct}
For all $r$ and $w$: $w \models r$ iff $\nu(\partial_w(r)) = \mathtt{True}$.
\end{lemma}
\begin{proof}
By induction on $w$, the base case being the definition of $\nu$, and the
step using derivative correctness for regular languages
\cite{brzozowski1964}, extended to \lstinline{Not} by
Equation~\eqref{eq:deriv-not} and to \lstinline{And} by the direct
derivative clause.
\end{proof}

Throughout this subsection we assume the single-word condition of
Theorem~\ref{thm:single-word}, without which \S\ref{sec:boundary} showed
that associativity holds only as an inclusion, so ``the residual of $e$''
would not be determined by $e$ alone and there would be nothing
well-defined to prove sound.
Under the assumption, a program $e$ built from primitives $p_1, \dots, p_n$
by \lstinline{pure} and \lstinline{>>=} emits one determinate trace
$w = u_1 \cdots u_n$, where $\post(p_i) = \{u_i\}$; write $w_{<i}$ and
$w_{>i}$ for the trace before and after $p_i$ runs.

\begin{theorem}[Soundness of $\RE$-residual checking]
\label{thm:soundness}
Let $e$ be a \lstinline{Pledge m (RE t) a} term built from \lstinline{pure}
and \lstinline{>>=} over primitives $p_1, \dots, p_n$ whose postconditions
each denote a single word.
If $L(\fut(e)) = \Sigma^*$ and $L(\pre(e)) = \Sigma^*$, then for every $i$:
\[
  w_{>i} \models \fut(p_i)
  \qquad\text{and}\qquad
  w_{<i} \models \pre(p_i).
\]
\end{theorem}
\begin{proof}
By induction on $n$; the base case $n = 0$ is vacuous.
For the step write $e = p \mathbin{>\!\!>\!\!=} f$ with $p = p_1$ and $f$
yielding $e'$, built from $p_2, \dots, p_n$, noting
$\post(e') = \{w_{>1}\}$.

\emph{Futures.} By Equation~\eqref{eq:fut-bind},
$\fut(e) = (\fut(p) \lquot \post(e')) \sqcap \fut(e')$; since $\sqcap$ is
intersection and $L(\fut(e)) = \Sigma^*$, both conjuncts denote $\Sigma^*$.
The first gives $L(\fut(p) \lquot \{w_{>1}\}) = \Sigma^*$, i.e.\
$w_{>1}\,v \models \fut(p)$ for every $v$; taking $v = \varepsilon$ gives
the claim for $i = 1$.
The second gives $L(\fut(e')) = \Sigma^*$, so the induction hypothesis
applies to $e'$, giving the claim for $i \geq 2$.

\emph{Preconditions.} Symmetrically,
$\pre(e) = \pre(p) \sqcap (\pre(e') \rquot \post(p))$, so both conjuncts
denote $\Sigma^*$; the first makes $w_{<1} = \varepsilon$ satisfy
$\pre(p)$, and the second, with $v = \varepsilon$, makes $u_1$ a trace
establishing $\pre(e')$, so the induction hypothesis on $e'$ gives the
claim for $i \geq 2$.
\end{proof}

\emph{It is a soundness result only}: a discharged residual certifies that
the obligations are met, but a non-discharged residual does not by itself
certify that they are not, which is exactly what
\S\ref{sec:pre-condition}'s past-closure condition protects.
\emph{The implementation tests a sufficient syntactic condition}: the
theorem is stated semantically ($L(\cdot) = \Sigma^*$), while
\lstinline{Pledge} computes $\norm{\cdot}$, a normaliser, and reports
discharge when the result is syntactically $\univ$ --- sound but incomplete
for language universality, so the check can fail to recognise a discharged
residual, never the reverse.
Theorem~\ref{thm:soundness} is the foundation both front-ends built in
\S\ref{sec:ltl}--\S\ref{sec:mtl} inherit unchanged, since both compile to
$\RE$ rather than reimplementing quotient or derivative.

\paragraph{A residual scope note.}
The \lstinline{Event} type of \S\ref{sec:re} still has no pattern between
an exact \lstinline{Atom} match and the catch-all \lstinline{Wildcard} ---
there is no way to say ``a \lstinline{write} with any arguments,'' the
first thing a specification author reaches for.
Adding a name-indexed wildcard would remove the restriction without
disturbing the algebra above: only \lstinline{subsumesEvent} and the
derivative's \lstinline{Single} clause depend on the exact-match reading.
We leave it as a small, self-contained item, not a limitation the theory
above depends on.

%% =========================================================================
\section{Runtime Monitoring with LTL, on Top of RE}
\label{sec:ltl}
%% =========================================================================

\S\ref{sec:effects} showed that a \lstinline{Pledge} annotation can
reference a real runtime value, and \S\ref{sec:soundness} proved that its
residual is sound with respect to the trace an execution actually emits.
Put together, these two facts license a reading of \lstinline{Pledge} that
is stronger than ``a specification'': because bind runs the first action
before computing the second, the residual after each step \emph{is} the
state of a monitor over the one execution that happened, and the quotient
\emph{is} the monitor's transition function.
This section makes that reading concrete with a proper linear-temporal-logic
front-end, \lstinline{Pledge.LTL}, whose formulas compile to the $\RE$ of
\S\ref{sec:re-instance} and whose monad laws and soundness theorem it
therefore inherits without further proof.

\subsection{The \texttt{LTL} Type}

LTL over finite traces~\cite{degiacomo2013} is the natural logic for the
obligations of \S\ref{sec:why}: a mutex is released, or a handle closed,
within the run, not somewhere along an infinite one.
\begin{lstlisting}
data LTL t
    = LTLTrue  | LTLFalse | LTLAtom (Event t) | LTLNot (LTL t)
    | LTLAnd (LTL t) (LTL t) | LTLOr (LTL t) (LTL t)
    | LTLNext (LTL t)          -- X phi
    | LTLUntil (LTL t) (LTL t) -- phi U psi
    | LTLFinally (LTL t)       -- F phi
    | LTLGlobally (LTL t)      -- G phi
    | LTLFromRE (RE t)         -- escape hatch: an opaque RE
\end{lstlisting}
Its standard translation goes through an automaton, but because complement
is a first-class $\RE$ constructor (\S\ref{sec:re}), every LTL operator
translates directly to $\RE$ with no automaton construction, via
\lstinline{ltlToRe} (Table~\ref{tab:ltl}); a best-effort inverse
\lstinline{reToLtl} recognises the $\RE$ shapes \lstinline{ltlToRe}
produces and is used only for pretty-printing.

\begin{table}[ht]
\caption{LTL to RE translation ($\Sigma^1 \triangleq \mathtt{Single\;Wildcard}$,
         $\univ \triangleq \neg\varnothing$).}
\label{tab:ltl}
\begin{tabular}{ll}
\toprule
LTL formula & RE $\sem{\cdot}$ \\
\midrule
$\top$              & $\univ$ \\
$\bot$              & $\varnothing$ \\
$e$                 & $\{e\}$ \\
$\neg\phi$          & $\neg\sem{\phi}$ \\
$\phi \wedge \psi$  & $\sem{\phi} \cap \sem{\psi}$ \\
$\phi \vee \psi$    & $\sem{\phi} \cup \sem{\psi}$ \\
$X\,\phi$           & $\Sigma^1 \cdot \sem{\phi}$ \\
$F\,\phi$           & $\univ \cdot \sem{\phi}$ \\
$G\,\phi$           & $\neg(\univ \cdot \neg\sem{\phi})$ \\
$\phi\,U\,\psi$     & $\mathit{step}(\phi)^* \cdot \sem{\psi}$
                      \quad ($\phi$ propositional) \\
\bottomrule
\end{tabular}
\end{table}

The \lstinline{Composable LTL} instance routes every operation through
\lstinline{ltlToRe}, wrapping the $\RE$ result back up as
\lstinline{LTLFromRE}, so \lstinline{concatenation}, the two quotients, and
the two identities are exactly $\RE$'s, viewed through the LTL syntax:
\begin{lstlisting}
instance Eq t => Composable (LTL t) where
  concatenation l1 l2 = LTLFromRE (concatenation (ltlToRe l1) (ltlToRe l2))
  conjunction         = LTLAnd
  empty               = LTLFromRE empty
  universe            = LTLTrue
  leftQuotient  l1 l2 = LTLFromRE (leftQuotient  (ltlToRe l1) (ltlToRe l2))
  rightQuotient l1 l2 = LTLFromRE (rightQuotient (ltlToRe l1) (ltlToRe l2))
\end{lstlisting}

\subsection{Monitoring a Mutex and a Transaction}
\label{sec:ltl-examples}

We read two of the scenarios from \S\ref{sec:why} as monitor runs: each
primitive posts an event and a future obligation, \lstinline{inspect}
executes the program exactly once, and the residual reports the monitor's
state at the end of that one execution.

\begin{lstlisting}
acquire :: Int -> Pledge IO (RE Values) ()
acquire mid = Pledge $ return
    ( ()
    , universe
    , Single (Atom "acquire" (List [Num mid]))
    , finally (Atom "release" (List [Num mid])) )

release :: Int -> Pledge IO (RE Values) ()
release mid = Pledge $ return
    ( ()
    , previously (Atom "acquire" (List [Num mid]))  -- whole-past precondition
    , Single (Atom "release" (List [Num mid]))
    , universe )
\end{lstlisting}
The precondition uses \lstinline{previously}, not \lstinline{Single}: by
\S\ref{sec:pre-condition} it must describe the whole preceding trace, or
\lstinline{nestedLocks} below --- where \lstinline{release 2} intervenes
between the acquire and its matching release --- would report a violation
for a correct run.

\begin{center}
\begin{tabular}{lll}
\toprule
Monitored run & Residual $\fut$ & Residual $\pre$ \\
\midrule
\lstinline{do acquire 1; release 1}                     & $\univ$ & $\univ$ \\
\lstinline{do acquire 1; acquire 2; release 2; release 1} & $\univ$ & $\univ$ \\
\lstinline{do acquire 1; acquire 2; release 2} (lock leak) & $\finally{\mathtt{release}(1)}$ & $\univ$ \\
\lstinline{release 1} (no matching acquire)              & $\univ$ & $\finally{\mathtt{acquire}(1)}$ \\
\bottomrule
\end{tabular}
\end{center}

\noindent
The lock-leak run's residual future names lock~1 specifically: lock~2's
obligation is discharged by \lstinline{release 2}'s postcondition, and
lock~1's $\finally{\mathtt{release}(1)}$ passes through the left quotient
unchanged.
The unmatched \lstinline{release 1}'s residual precondition
$\finally{\mathtt{acquire}(1)}$ reads as an obligation transferred to the
caller: composing this fragment after \lstinline{acquire 1} discharges it
to $\univ$.
This is the case a lexically scoped combinator such as \lstinline{bracket}
handles least well: two obligations are outstanding when
\lstinline{release 2} runs, and which one it discharges is not determined
by the program's nesting, only by which event actually occurred --- exactly
what a monitor, and not a static scope, can see.

A disjunctive obligation makes the same point from the other direction.
\lstinline{beginTx} posts a future requiring eventual \lstinline{commit} or
\lstinline{rollback}; whichever event the run actually emits discharges the
whole disjunction via the left quotient, and \lstinline{beginTx} alone
retains it in full.
No pre/post pair can express ``resolved by either of two later events, one
of which is chosen at runtime''; a monitor whose state is the residual
itself can, because that state is exactly a set of continuations, not a
Boolean flag.

%% =========================================================================
\section{Static Verification with GuardedRE, on Top of RE}
\label{sec:gre}
%% =========================================================================

$\RE$ and its $\LTL$ front-end are silent about arithmetic state, and their
soundness theorem (\S\ref{sec:soundness}) is inherently about \emph{one}
executed trace: checking it is runtime monitoring.
The $\GRE$ instance adds a second, independent dimension --- a Presburger
arithmetic predicate over heap addresses, discharged by Z3 via
\lstinline{sbv} --- and because that predicate is checked over
\emph{symbolic} heap values rather than the concrete heap of one run, the
same query can be asked \emph{before} any execution: is this residual
guard satisfiable at all, for some heap; is it valid for every heap
satisfying an invariant.
That is what makes $\GRE$ a static-verification instance rather than a
second runtime monitor.

\subsection{The \texttt{GuardedRE} Type}

\begin{lstlisting}
type GuardedRE a = [(Pred, RE a)]
\end{lstlisting}
A state $(\text{heap}, \text{trace})$ satisfies \lstinline{gre :: GuardedRE a}
iff \emph{some} disjunct $(p, r) \in \mathit{gre}$ has
$\text{heap} \models p$ and $\text{trace} \in L(r)$.
Within one disjunct the two dimensions are independent --- the predicate is
static and the RE is advanced by derivatives --- but the list carries
\emph{disjunction} across dimensions: a plain conjunctive pair $(p, r)$
cannot express ``the heap is in state $A$ and the trace obeys $r_A$, or the
heap is in state $B$ and the trace obeys $r_B$'', which is exactly the
shape a property that changes as the heap evolves takes, and is what an
earlier, single-conjunct design could not express.
Smart constructors lift a pure RE or a pure predicate:
\begin{lstlisting}
fromRE    :: RE a  -> GuardedRE a
fromPPred :: Pred  -> GuardedRE a
\end{lstlisting}

\subsection{The \texttt{Composable GuardedRE} Instance}

\begin{lstlisting}
instance Eq a => Composable (GuardedRE a) where
  concatenation xs ys = [ (normalizePred (PAnd p1 p2), normalize (Seq r1 r2))
                         | (p1,r1) <- xs, (p2,r2) <- ys ]
  conjunction   xs ys = [ (normalizePred (PAnd p1 p2), normalize (And r1 r2))
                         | (p1,r1) <- xs, (p2,r2) <- ys ]
  empty    = [(PTrue, Epsilon)]
  universe = [(PTrue, top)]
  leftQuotient  xs ys = [ (normalizePred (PAnd p1 p2), normalize (reLeftQuotient r1 r2))
                         | (p1,r1) <- xs, (p2,r2) <- ys ]
  rightQuotient xs ys = [ (normalizePred (PAnd p1 p2), normalize (reRightQuotient r1 r2))
                         | (p1,r1) <- xs, (p2,r2) <- ys ]
\end{lstlisting}
Every operation is the cross product of disjuncts, conjoining the Pred
halves and applying the corresponding $\RE$ operation to the trace halves;
a normaliser (\lstinline{normalizeGuarded}) drops unsatisfiable-looking
disjuncts and merges disjuncts that share a predicate, keeping the list
from growing unboundedly under repeated composition, and an SMT-backed
variant (\lstinline{normalizeGuardedSMT}) additionally collapses disjuncts
whose predicates are jointly exhaustive --- e.g.\ a two-way heap split
$h[a] = 0 \vee h[a] > 0$ collapses to a single \lstinline{PTrue} disjunct
--- which needs a solver call to see and so lives in \lstinline{IO}.

\subsection{Presburger Predicates, and Two Kinds of Solver Query}

\begin{lstlisting}
data Term = Val Values | Add Term Term | Neg Term | Mul Int Term
data Pred = PTrue | PFalse
          | PLt Term Term | PLe Term Term | PEq Term Term
          | PGt Term Term | PGe Term Term
          | PNot Pred | PAnd Pred Pred | POr Pred Pred
\end{lstlisting}
\lstinline{Val (ValAt a)} reads the value at heap address \lstinline{a}.
Two solver queries give the static reading its two flavours:
\lstinline{checkPPred :: Pred -> IO SolverResult} asks Z3 for a witness heap
satisfying a predicate (an existential, SAT query over free heap and
value variables), and \lstinline{isValidUnderHeapInvariant :: Pred -> IO Bool}
asks whether a predicate holds for \emph{every} heap obeying the standing
invariant that addresses are non-negative, by checking that its negation,
conjoined with that invariant, is unsatisfiable.
Neither call needs a concrete run: both are posed over the symbolic
addresses the predicate mentions.
\lstinline{checkGuarded} folds \lstinline{deriveGuarded} (which advances
only the RE half of every disjunct) over a trace and then asks
\lstinline{checkPPred} on the nullable disjuncts against a \emph{concrete}
heap, giving $\GRE$ a runtime reading as well when one is wanted --- the two
modes share the same representation and differ only in whether the
predicate side is instantiated or left symbolic.

\subsection{Heap Liveness and Allocation Order}
\label{sec:ex-gre-memory}

Manual memory management needs a trace constraint (\lstinline{malloc} must
precede its matching \lstinline{free}, which must eventually happen, and a
second \lstinline{free} must not occur before re-allocation) and a heap
constraint (the cell is live, $h[a] > 0$) at once; an RE alone cannot see
liveness, and a predicate alone cannot see ordering.
Each \lstinline{malloc}/\lstinline{free} pair contributes \emph{two}
disjuncts, one per heap state the cell can be in, so the obligation tracks
a property that changes as the trace proceeds rather than only a static
invariant:
\begin{lstlisting}
malloc :: Pledge IO (GuardedRE Values) Addr
malloc = Pledge $ do
    addr <- randomRIO (0, 5)
    return ( addr
           , fromRE universe
           , [ (PEq (Val (ValAt addr)) (Val (Num 0)), Epsilon)
             , (PGt (Val (ValAt addr)) (Val (Num 0)), Single (Atom "malloc" (List [Num addr]))) ]
           , [ (PEq (Val (ValAt addr)) (Val (Num 0)), never (Usage (List [Num addr])))
             , (PGt (Val (ValAt addr)) (Val (Num 0)), finally (Atom "free" (List [Num addr]))) ]
           )

free :: Addr -> Pledge IO (GuardedRE Values) ()
free addr = Pledge $ return
    ( ()
    , [ (PEq (Val (ValAt addr)) (Val (Num 0)), universe)
      , (PGt (Val (ValAt addr)) (Val (Num 0)), previously (Atom "malloc" (List [Num addr]))) ]
    , fromRE (Single (Atom "free" (List [Num addr])))
    , fromRE (noUntil (Atom "free" (List [Num addr])) (Atom "malloc" (List [Num addr]))) )
\end{lstlisting}

\begin{center}
\small
\begin{tabular}{lll}
\toprule
Program & Residual $\pre$ & Residual $\fut$ \\
\midrule
\lstinline{mallocThenFree}, \lstinline{reallocate} (good)
  & $[h{=}0]\,\univ \lor [h{>}0]\,\univ$
  & $[h{>}0]\,\noUntil{\mathtt{free}}{\mathtt{malloc}}$ \\
\lstinline{missingFree} (leak)
  & $[\mathbf{true}]\,\univ$
  & $[h{>}0]\,\finally{\mathtt{free}}$ \\
\lstinline{freeWithoutMalloc}
  & $[h{>}0]\,\finally{\mathtt{malloc}}$
  & --- \\
\lstinline{doubleFree}
  & $[h{=}0]\,\univ \lor [h{>}0]\,\univ$
  & $\varnothing$ \\
\bottomrule
\end{tabular}
\end{center}

\noindent
\lstinline{doubleFree} reaches an empty future because the second
\lstinline{free} must satisfy the \lstinline{noUntil} guard the first one
imposed, and no trace does; \lstinline{reallocate} shows the same guard
reset cleanly by an intervening \lstinline{malloc}.
Between $\univ$ and $\varnothing$ lies a satisfiable-but-not-universal
residual, which reads as a standing constraint on whatever runs next rather
than a verdict --- \lstinline{mallocThenFree} is correct and still carries
$\noUntil{\mathtt{free}}{\mathtt{malloc}}$ forward, precisely because that
obligation belongs to the continuation, not to the pair itself.

%% =========================================================================
\section{Checking Out MTL, on Top of RE}
\label{sec:mtl}
%% =========================================================================

$\LTL$ (\S\ref{sec:ltl}) is qualitative about time: $F\,\phi$ says $\phi$
happens \emph{eventually}, with no bound on when.
Real protocols are often quantitative instead --- a lock must be released
\emph{within five steps}, a valve must close \emph{between two and four
steps} after opening.
\lstinline{Pledge.MTL} adds bounded, discrete-step intervals to the same
syntax, and --- like $\LTL$ --- compiles down to $\RE$ rather than adding a
new algebra, so it inherits \S\ref{sec:re-instance}'s monad laws and
soundness theorem for free.

\subsection{Bounded Intervals and the \texttt{MTL} Type}

\begin{lstlisting}
data Interval = Interval Int (Maybe Int)   -- [lo, hi] or [lo, infinity)

within, exactly, atLeast :: Int -> Interval
fromTo :: Int -> Int -> Interval

data MTL t
    = MTLTrue | MTLFalse | MTLAtom (Event t) | MTLNot (MTL t)
    | MTLAnd (MTL t) (MTL t) | MTLOr (MTL t) (MTL t) | MTLNext (MTL t)
    | MTLUntil    Interval (MTL t) (MTL t)  -- phi U_[lo,hi] psi
    | MTLFinally  Interval (MTL t)          -- F_[lo,hi] phi
    | MTLGlobally Interval (MTL t)          -- G_[lo,hi] phi
    | MTLFromRE   (RE t)
\end{lstlisting}
A bounded modality compiles by finite unrolling: $F_{[lo,hi]}\,\phi$ becomes
a union, over each step count $k \in [lo, hi]$, of ``skip exactly $k$
events, then $\phi$'':
\begin{lstlisting}
mtlToRe (MTLFinally i p) = do
    ks <- intervalSteps i
    r  <- mtlToRe p
    pure $ case ks of
        Left  a  -> Seq (sigmaPow a) (Seq top r)          -- unbounded above: [a, infinity)
        Right xs -> foldr Or Bot [ Seq (sigmaPow k) r | k <- xs ]  -- [lo, hi]
\end{lstlisting}
where \lstinline{sigmaPow n = Seq (Single Wildcard) (sigmaPow (n-1))}
concatenates exactly $n$ wildcards.
$G_{[lo,hi]}\,\phi$ and $\phi\;U_{[lo,hi]}\,\psi$ unroll the same way,
using \lstinline{atMostSeq} and \lstinline{reptSeq} respectively to bound
the number of repetitions; a best-effort inverse \lstinline{reToMtl}
mirrors \S\ref{sec:ltl}'s \lstinline{reToLtl} for pretty-printing.
The \lstinline{Composable MTL} instance is the same shape as
\lstinline{Composable LTL}: every operation round-trips through
\lstinline{mtlToRe} and wraps the $\RE$ result as \lstinline{MTLFromRE}, so
an $\MTL$ specification is, after compilation, an ordinary $\RE$
specification that \S\ref{sec:re-instance} already knows how to discharge.

\subsection{A Bounded Obligation}

Three combinators cover the common patterns directly, without spelling out
an \lstinline{Interval} at each use site:
\begin{lstlisting}
finallyWithin :: Int -> Event t -> MTL t   -- F_[0,n] ev
neverWithin   :: Int -> Event t -> MTL t   -- not F_[0,n-1] ev
globallyWithin:: Int -> Event t -> MTL t   -- G_[0,n-1] ev
\end{lstlisting}
Annotating a valve's \lstinline{open} with
\lstinline{finallyWithin 5 (Atom "close" v)} as its future obligation reads,
after compilation, as a five-way union of ``close occurs after exactly $k$
other events'' for $k \in [0,5]$; a run that closes the valve on step~$3$
discharges it via the same left-quotient machinery as an unbounded
\lstinline{finally}, and a run that never closes it leaves the same kind of
non-$\univ$ residual as \S\ref{sec:ltl-examples}'s lock leak, now additionally
carrying the deadline that was missed.
A worked case study exercising \lstinline{MTL} end to end, alongside the
\S\ref{sec:ltl-examples} and \S\ref{sec:ex-gre-memory} examples, is not yet
in the \texttt{Examples} tree; adding one is a short, self-contained task
we list under \S\ref{sec:experiments}.

%% =========================================================================
\section{Usability: Making Error Location Traceable}
\label{sec:usability}
%% =========================================================================

Everything up to this point reports a violation as a value: a non-$\univ$
$\fut$ or a collapsed $\pre$.
Nothing yet says \emph{where} in a composed program that value first
stopped being $\univ$.
For a two-line example this does not matter; for the kind of program
\lstinline{Pledge} is meant to specify, it does, so we record this as a
concrete, currently unimplemented piece of the design rather than fold it
into the discussion.

\paragraph{The problem.}
\lstinline{inspect} runs the whole composed action and reports only the
final residual; a violation folded in at bind step 3 of 30 looks identical
to one folded in at step 29.
The programmer must bisect the \lstinline{do}-block by hand to find the
offending primitive.

\paragraph{Plan.}
\begin{itemize}[wide, labelindent=10pt]
  \item Add a labelling combinator,
        \lstinline{label :: String -> Pledge m eff a -> Pledge m eff a}
        (or a \lstinline{HasCallStack} variant that needs no explicit
        string), that tags one primitive's call site without changing its
        annotations.
  \item Thread an accumulating trace of \lstinline{(label, pre, post, fut)}
        triples through bind, alongside the existing 4-tuple, so a
        labelled step's contribution to the residual is individually
        recoverable rather than only visible in the folded-in result.
  \item Report, on inspection, the \emph{first} labelled step at which
        $\pre$ became $\varnothing$ or the running $\fut$ stopped being
        $\univ$ --- turning ``the program violates its contract'' into
        ``line \emph{N} does''.
  \item Make this opt-in and free when unused: a program that never calls
        \lstinline{label} pays no overhead beyond what \S\ref{sec:pledge}
        already costs.
\end{itemize}
This closes a gap independently noted in early user feedback on the
library (violations are silent and unlocated) and pairs naturally with a
second, related item: nothing currently stops execution when $\pre$
collapses to $\varnothing$, so a \lstinline{runChecked} that inspects the
residual after each bind and fails as soon as it reaches $\varnothing$
would turn a located violation into an enforced one, at the cost of a
quotient computation per bind.
Both are additions to \lstinline{Pledge.Core}; neither requires revisiting
\S\ref{sec:monad-laws}--\S\ref{sec:mtl}.

%% =========================================================================
\section{Auto-Generation of Proofs}
\label{sec:proof-gen}
%% =========================================================================

\S\ref{sec:monad-laws} proves the monad laws once, abstractly, from twelve
axioms; \S\ref{sec:mechanisation} then discharges those twelve axioms for
$\RE$ by hand in Lean.
Every new \lstinline{Composable} instance --- $\GRE$ today, whatever comes
after it tomorrow --- currently repeats that second step from scratch: a
person mirrors the Haskell definition in Lean and proves twelve lemmas
again.
This does not scale past a handful of instances, and it is the piece of the
mechanisation story we have not automated.

\paragraph{Plan.}
\begin{itemize}[wide, labelindent=10pt]
  \item \textbf{Property-test first.} Before writing any Lean, state the
        twelve axioms as QuickCheck properties over the candidate instance
        --- exactly what \S\ref{sec:mechanisation} already does for $\RE$
        --- and run them over a bounded pool. A property that fails here
        needn't be attempted in Lean; \S\ref{sec:boundary} is what a
        systematic sweep of this kind found for $\RE$.
  \item \textbf{Generate the Lean statements, not the proofs, mechanically.}
        The twelve axiom statements are the same schema for every instance,
        parametrised only by the instance's \lstinline{concatenation},
        \lstinline{conjunction}, and quotients; a small script can emit the
        Lean lemma \emph{signatures} for a new instance directly from its
        Haskell type-class instance, leaving only the proof bodies to fill
        in.
  \item \textbf{Discharge the finite fragment automatically.} Several
        axioms (S1--S3, C1--C3 in Table~\ref{tab:axioms}) are generic
        properties of an underlying monoid/semilattice pair and are
        provable by \lstinline{decide} or a Mathlib tactic once the
        instance's carrier is finite or the operations are given as
        pattern matches, without a hand-written proof term; only L2--L4 and
        R2--R4 need instance-specific reasoning about the quotient.
  \item \textbf{Track the boundary as data.} Record, per instance, exactly
        which axioms hold unconditionally, which hold under a stated side
        condition, and which fail --- as Table~\ref{tab:axiom-status} does
        for $\RE$ --- so that adding an instance produces a table entry
        automatically alongside whatever manual proof the residual cases
        still need.
\end{itemize}
The target is not a fully automatic prover for arbitrary instances, which
is out of scope; it is removing the boilerplate that is already mechanical
today, so that mechanising a fifth instance costs a fraction of what the
first two did.

%% =========================================================================
\section{Experiments}
\label{sec:experiments}
%% =========================================================================

The evaluation so far is entirely about correctness --- the property-based
validation of \S\ref{sec:mechanisation} and the case studies of
\S\ref{sec:ltl-examples} and \S\ref{sec:ex-gre-memory} --- and says nothing
about cost.
That evaluation is not yet done; we scope it concretely rather than claim
results we do not have.

\paragraph{Plan.}
\begin{itemize}[wide, labelindent=10pt]
  \item \textbf{Quotient and residual-size growth.} For the $\RE$ case
        studies of \S\ref{sec:ltl-examples}, measure derivative/quotient
        wall-clock time and normalised-residual size as a function of trace
        length, to check that the worklist termination argument of
        \S\ref{sec:composable-re} stays cheap in practice and not only in
        the worst case.
  \item \textbf{SMT call cost for $\GRE$.} For \S\ref{sec:ex-gre-memory},
        count and time the \lstinline{checkPPred} /
        \lstinline{isValidUnderHeapInvariant} calls per program, and report
        how much \lstinline{normalizeGuardedSMT} reduces disjunct-list
        growth compared to the purely structural
        \lstinline{normalizeGuarded}.
  \item \textbf{Bounded vs.\ unbounded compilation.} Compare $\MTL$
        (\S\ref{sec:mtl}) against $\LTL$ (\S\ref{sec:ltl}) on the same
        obligation as the interval width grows, to characterise the
        unrolling blow-up \lstinline{sigmaPow}/\lstinline{reptSeq}
        introduce and find where it stops being negligible.
  \item \textbf{Scale up the axiom validation.} Extend
        \S\ref{sec:mechanisation}'s bounded-pool QuickCheck sweep with
        larger pools and longer words, and add the $\GRE$ instance to it,
        both as a correctness check and as the first input to
        \S\ref{sec:proof-gen}'s property-test-first step.
  \item \textbf{Report, don't assert.} Present the above as tables/plots
        over the existing \texttt{Examples} programs plus the one new $\MTL$
        case study noted in \S\ref{sec:mtl}, rather than a synthetic
        benchmark suite, so the numbers describe the same programs the rest
        of the paper already discusses.
\end{itemize}

%% =========================================================================
\section{Related Work}
\label{sec:related}
%% =========================================================================

\paragraph{Future conditions and temporal verification.}
The proposal this paper implements is due to \citet{song2025futurecond},
who introduce future conditions as a first-class extension to pre/post
contracts at the level of a formal system.
\citet{song2024provenfix} apply a related RE encoding of temporal
properties in ProveNFix to guide automated program repair.
Pledge's contribution is a concrete, runnable library realisation: a monad
transformer, a derivative-based discharge mechanism, and a mechanised
account of exactly when it is lawful, neither prior paper gives.

\paragraph{Pre/post-condition contracts in Haskell.}
Liquid Haskell~\cite{vazou2014} expresses pre/post-conditions as refinement
types verified by an SMT solver.
Findler and Felleisen~\cite{findler2002} establish runtime contract discipline
for higher-order functions.
Pledge complements these by adding a temporal future condition without
requiring type-system extensions.

\paragraph{Specifications as monads.}
The idea that a specification can itself be monadic structure is the
closest conceptual neighbour to this work.
Hoare Type Theory~\cite{nanevski2008htt} indexes a monad by a pre/post
pair, making Hoare triples types.
The \emph{Dijkstra monad}~\cite{swamy2013dijkstra} instead indexes
computations by predicate transformers, so that the weakest precondition of
a composite is computed by bind;
\citet{ahman2017dijkstra} derive such monads mechanically from an effect's
definition, and \citet{maillard2019dijkstra} give the general account in
which a Dijkstra monad is a monad morphism into a specification monad.
\citet{swierstra2019} develop the predicate-transformer view for effects in
a dependently typed setting.

The shared idea is that specifications compose the way programs do; the
difference is what composition computes.
A Dijkstra monad propagates \emph{backwards}: bind computes the weakest
precondition of a sequence from its continuation's, and the result is a
predicate transformer to be discharged by a verifier.
\lstinline{Pledge} propagates \emph{forwards} as well, and its $\fut$ field
has no counterpart in that line of work --- a predicate transformer relates
two states, whereas a future condition constrains the trace still to come,
which no pre/post pair over states can express.
The second difference is one of ambition.
Dijkstra monads target full functional correctness in a dependently typed
host and rely on a solver; \lstinline{Pledge} targets temporal obligations
in plain Haskell, computes a residual rather than a verification condition,
and is correspondingly weaker and cheaper.

\paragraph{Parameterised and graded monads.}
\citet{atkey2009} introduces parameterised monads for Hoare-style state
tracking; \citet{katsumata2014} generalises this to graded monads;
\citet{orchard2016} shows effect systems and session types are instances.
Pledge carries obligations as value-level annotation fields rather than
indexing by effect type at the type level, preserving compatibility with the
standard \lstinline{Monad} class and \lstinline{do}-notation.

\paragraph{Resource management in Haskell.}
The \lstinline{bracket} combinator lexically scopes cleanup.
\lstinline{ResourceT}~\cite{snoyman2011resourcet}, \lstinline{Acquire}, and
\lstinline{Managed} thread finaliser queues through every bind.
All three handle the cleanup half of a resource obligation but cannot express
protocol obligations, which require explicit future-condition propagation
across bind steps, as \S\ref{sec:ltl-examples}'s transaction example shows.

\paragraph{Linear types.}
Linear Haskell~\cite{bernardy2018linear} enforces ``consume exactly once''.
Linear types enforce structural obligations (use at most once), whereas
future conditions enforce temporal obligations (use in a specified order,
or within a specified deadline as in \S\ref{sec:mtl}).
A hybrid combining both could achieve stronger guarantees than either alone.

\paragraph{Session types.}
Session types~\cite{honda1993,gay2010} enforce communication protocols at the
type level.
The \lstinline{sessiontypes} library~\cite{sessiontypes2017} realises this
via an indexed monad.
Pledge is more lightweight: it operates within the standard \lstinline{Monad}
class and is applicable beyond message-passing.

\paragraph{Typestate.}
Typestate~\cite{strom1986} attaches a protocol state to a variable's type
and checks transitions statically; \citet{bierhoff2007} extend this to
aliased objects with access permissions.
The obligations are the same ones \lstinline{Pledge} targets --- an opened
file must be closed, a lock released --- but the mechanism differs in where
the state lives.
Typestate puts it in the type, so the checker rejects a bad program and the
protocol must be fixed before the code is written;
\lstinline{Pledge} puts it in a value, so a bad program compiles and its
residual reports what it owes, located traceably once \S\ref{sec:usability}
is in place.

\paragraph{Runtime verification.}
Runtime verification~\cite{leucker2009} checks temporal properties against a
concrete execution rather than a model, and is the setting
\S\ref{sec:ltl} argues \lstinline{Pledge} realises.
\citet{havelund2001rewriting} monitor temporal formulae by rewriting them as
events arrive --- the residual formula after each event being precisely the
obligation that remains, the same move the quotient makes here.
MOP~\cite{chen2007mop} generates efficient monitors from a specification
attached to program points, and Copilot~\cite{pike2010copilot} generates
hard-real-time C monitors from a Haskell EDSL, the closest prior system to
\S\ref{sec:mtl}'s bounded intervals.
Two things distinguish the present work.
First, most monitors are automaton- or rewriting-based and report a Boolean
verdict; here the monitor state \emph{is} the specification, in the same
algebra the programmer wrote, so a failure names the outstanding obligation.
Second, monitor composition is usually external --- monitors are attached to
a program, and combining them is a matter of running several at once ---
whereas here composition is the program's own \lstinline{>>=}, governed by
the algebraic laws of \S\ref{sec:monad-laws}.
The cost, as \S\ref{sec:why} states plainly, is that \lstinline{Pledge}
observes only what the annotations report and does not yet intervene ---
closing that gap is \S\ref{sec:usability}'s enforcing-mode item.

\paragraph{Derivatives of regular expressions.}
Brzozowski derivatives~\cite{brzozowski1964} have found application in
parsing~\cite{owens2009,might2011} and model checking;
\citet{antimirov1996} refine them into partial derivatives, whose sets stay
smaller.
Equation~\eqref{eq:deriv-not} extends the construction to complement.
What is new here is not the derivative but the use of a \emph{language}
quotient --- of one specification by another, rather than by a single event
--- as the propagation rule of a monadic bind, and the two front-ends of
\S\ref{sec:ltl}--\S\ref{sec:mtl} that reuse it unchanged.

%% =========================================================================
\section{Conclusion}
\label{sec:conclusion}
%% =========================================================================

We have presented \emph{Pledge}, a Haskell library that realises
\citet{song2025futurecond}'s proposal of future conditions as a concrete
monad transformer: a \lstinline{Pledge m eff a} computation in underlying
monad \lstinline{m} that carries three annotation fields
$(\pre, \post, \fut) :: \mathit{eff}$, propagated algebraically through
monadic bind via left- and right-quotient operations abstracted as a
six-operation \lstinline{Composable} class.

The design now spans two solid instances built on one theory.
$\RE$, with Brzozowski derivatives computing both quotients, satisfies
twelve algebraic axioms exactly under two stated side conditions ---
single-word postconditions, past-closed preconditions --- proved in Lean~4
alongside a soundness theorem connecting a discharged residual to the trace
a run actually emitted.
On top of $\RE$, an $\LTL$ front-end reframes \lstinline{Pledge} as a
compositional runtime monitor whose state is the residual itself, and a
bounded $\MTL$ front-end adds discrete-time deadlines by compiling to the
same $\RE$.
Beside it, a fourth, disjunctive instance $\GRE$ pairs traces with
Presburger heap predicates discharged by Z3, giving \lstinline{Pledge} a
genuinely static reading as well as a runtime one, checked over symbolic
rather than concrete heaps.

What the design does not yet do is the honest measure of what remains:
a violation is reported but not located
(\S\ref{sec:usability}); the Lean proof burden of a new instance is still
manual (\S\ref{sec:proof-gen}); and the design has been validated for
correctness but not yet measured for cost (\S\ref{sec:experiments}).
A few smaller items sit beside these three: promoting the residual from a
reported value to an enforced one by failing bind on $\varnothing$
(\S\ref{sec:usability}); a parallel composition operator, corresponding to
the shuffle product on $\RE$, for concurrent programs, which the
sequential \lstinline{Pledge} monad does not yet offer; and hardening the
API so that \S\ref{sec:single-word} and \S\ref{sec:pre-condition}'s side
conditions are unwriteable rather than merely documented, by exposing
smart constructors in place of the raw $\RE$ constructors in $\post$ and
$\pre$ slots.
None of these require revisiting the theory of \S\ref{sec:monad-laws}--\S\ref{sec:re-instance};
each is a self-contained piece of engineering on top of a foundation we
have proved correct.

%% =========================================================================
\bibliographystyle{ACM-Reference-Format}
\bibliography{references}

\end{document}
